\section{Averaging the unramified factors}\label{sec:sieve}

Proposition~\ref{prop:additive-equidistribution} distributes the
nonnegative mass $\phi(z)\prod_{p\in S}f_p^*(z)$ among residue classes at
primes outside $S$.  To prove Theorem~\ref{thm:exterior-count}, we must
insert the remaining
factors $b_p(z)$ while retaining their local means.  Some of those means are
less than one, so the argument must also keep track of the probability that
a prime does not divide any component of $z$.  We first obtain a weighted
sieve estimate with this feature and then combine its local means with the
Euler factors in Proposition~\ref{prop:ext-multiplicity-majorant}.

\subsection{Residue functions and the local bounds}

Fix $m\in\Z\setminus\{0\}$ and a bounded region of the real trace-$m$
plane.  Throughout the sieve argument, $\phi$ is one of the nonnegative
smooth box functions in Proposition~\ref{prop:additive-equidistribution},
with its planar support in this region.  Its two scales lie between
$X^{-\tau}$ and
fixed constants.  Define a measure on $F$ by
\begin{equation}\label{eq:sieve-measure}
 \nu_\phi
 =\sum_{\substack{z\in F,\ \Tr z=m\\z\text{ integral outside }S}}
       \phi(z)\prod_{p\in S}f_p^*(z)\,\delta_z,
 \qquad
 V_\phi=A_\phi\prod_{p\in S}d_p(m).
\end{equation}
Here $\delta_z$ is unit point mass at $z$.  All its coefficients are
nonnegative, and Proposition~\ref{prop:additive-equidistribution} with
modulus one gives a finite total mass.
Lemma~\ref{lem:fourier-scalar-density} gives
\begin{equation}\label{eq:sieve-volume}
 V_\phi\asymp A_\phi
\end{equation}
uniformly in $S$: for large primes the logarithms of $d_p(m)=1+O(p^{-2})$
have an absolutely convergent sum, and the finitely many small-prime
densities have uniform positive upper and lower bounds.

For a positive integer $l$ prime to every prime of $S$, the integral
trace-$m$ coset has $l^2$ residue classes modulo $l$ in its two integral
coordinates.  Write $\E_l W$ for the uniform mean of a function $W$ on
these classes, and also write $W(z)$ for its pullback to points integral
outside $S$.

\begin{lemma}[Nonnegative residue functions]\label{lem:sieve-residue-function}
With $\phi$ and $V_\phi$ as in Equation~\eqref{eq:sieve-measure}, let
$1\le l\le X^\theta$ be an integer prime to every prime of $S$.
For every nonnegative
function $W$ on the residue classes of the integral trace-$m$ coset
modulo $l$,
\begin{equation}\label{eq:sieve-residue-function}
 \int W\,d\nu_\phi
   =V_\phi\E_l W+O\bigl(V_\phi X^{-\theta}\E_l W\bigr).
\end{equation}
The implied constant is uniform in the fields, orders, embedding
orderings, permitted box locations and scales, $l$, and $W$.  It may
depend on $m$, the fixed bounded region, and the fixed smooth box function.
\end{lemma}

\begin{proof}
For each residue class $a$, Equation~\eqref{eq:add-equidistribution} gives
\[
 \nu_\phi(a)=V_\phi l^{-2}(1+\epsilon_a),
 \qquad |\epsilon_a|\le C_{\mathrm{add}}X^{-\theta},
\]
with the same $C_{\mathrm{add}}$ for every $a$ and $l$ in the stated range.
Multiply by
$W(a)$ and sum.  Since $W(a)\ge0$, the absolute value of the summed error
is at most
$C_{\mathrm{add}}V_\phi X^{-\theta}l^{-2}\sum_a W(a)
 =C_{\mathrm{add}}V_\phi X^{-\theta}\E_l W$.
This also covers $\E_l W=0$.  In particular the error retains the local
mean even when that mean is very small.
\end{proof}

We recall the information about the functions to be inserted.  For
$p\notin S$, let $\E_p$ and $\P_p$ denote expectation and probability for
normalized additive Haar measure on
$\{z\in R_{F,p}:\Tr z=m\}$.  On $R_{F,p}\cap F_p^\times$, recall that
$n_p(z)=\max_{\mathfrak q\mid p}v_{\mathfrak q}(z)$ is the largest
component valuation, with $v_{\mathfrak q}(p)=1$; see
Equation~\eqref{eq:ext-unramified-factor}.  On the null set where a
component vanishes, keep the convention $n_p=\infty$, $b_p=0$.
Lemma~\ref{lem:ext-unramified-data} permits one fixed integer $C\ge1$
such that, on the nonvanishing locus,
\begin{align}
 &b_p\ge0,\qquad b_p=1\ \text{if }n_p=0,\qquad
 b_p\le C(1+n_p)^C,\label{eq:sieve-local-weight}\\
 &\P_p(n_p\ge j)\le Cp^{-j}\quad(j\ge1),\qquad
 \vartheta_p:=\P_p(n_p\ge1)\le C/p,\quad
 \vartheta_p\le\tfrac12. \label{eq:sieve-local-tail}
\end{align}
The event $\{n_p\ge j\}$ and the function
$b_p\mathbf 1_{\{n_p=j\}}$ are determined by the two integral coordinates
modulo $p^{C(j+1)}$ for every integer $j\ge1$.  They have extensions
defined by congruence classes on the whole local coset, agreeing with their values at every
global point.  A prime with $n_p(z)\ge1$ will be called active at $z$.
In particular,
\begin{equation}\label{eq:sieve-mean-lower}
 \E_p b_p\ge1-\vartheta_p\ge\tfrac12.
\end{equation}

Lemma~\ref{lem:ext-support} places the support of $\nu_\phi$ in the
trace-$m$ coset of $(C_F^\sharp)^\vee$.  On this support all $n_p(z)$,
$p\notin S$, are finite.  The lemma's norm and depth conclusions bound
the product of the depth factors by $O(X^{84})$ in our fixed bounded real region.  Thus,
for a fixed $B>84$ and all sufficiently large $X$,
\begin{equation}\label{eq:sieve-depth-product}
 \prod_{p\notin S}p^{n_p(z)}\le X^B
 \qquad\bigl(\nu_\phi(\{z\})>0\bigr).
\end{equation}

Two consequences will be used repeatedly.  There is a fixed $A_0>1$ such
that $b_p(z)\le A_0^{n_p(z)}$.  This follows from
Equation~\eqref{eq:sieve-local-weight} for $n_p\ge1$, since
$\sup_{j\ge1}(C(1+j)^C)^{1/j}<\infty$, and from $b_p=1$ at depth zero.
Also, with
\[
 \mathcal W(z)=\prod_{p\notin S}b_p(z),
\]
the product is finite on the support of $\nu_\phi$ and satisfies
\begin{equation}\label{eq:sieve-pointwise-weight}
 \mathcal W(z)\ll_\epsilon X^\epsilon
 \qquad(\epsilon>0).
\end{equation}
Indeed, for any $t>0$, the inequality
$C(1+j)^C\le p^{tj}$ holds for all $j\ge1$ once $p$ is sufficiently large.
For each of the remaining finitely many primes it holds with a constant
factor independent of $j$.  Multiplying these inequalities and using
Equation~\eqref{eq:sieve-depth-product}, with $t=\epsilon/B$, proves
Equation~\eqref{eq:sieve-pointwise-weight}.  All constants here are uniform
in the arithmetic data and in the permitted boxes.

\begin{proposition}[Weighted averaging]\label{prop:sieve-averaging}
For the measure $\nu_\phi$ in Equation~\eqref{eq:sieve-measure}, there is
a fixed $\delta>0$ such that
\begin{equation}\label{eq:sieve-average}
 \int\mathcal W(z)\,d\nu_\phi(z)
   \ll V_\phi\prod_{\substack{p\le X^\delta\\p\notin S}}\E_p b_p.
\end{equation}
The choice of $\delta$ and the implied constant are independent of the
field, order, embedding ordering, and the permitted box location and
scales.  They may depend on $m$, the fixed bounded real region, and the
fixed smooth box function.
\end{proposition}

The proof first counts bounded products of prescribed positive depths,
while requiring inactivity at any other chosen small primes.  The
inactivity condition preserves the factors in
Equation~\eqref{eq:sieve-mean-lower}.  We then bound the contribution of two
exceptional sets of points and truncate the remaining depth products within
the modulus range of
Lemma~\ref{lem:sieve-residue-function}.

\subsection{Prescribed depths and inactive primes}

Write
$\mathcal P_y=\{p\le y:p\notin S\}$.  Chebyshev's estimate
$\sum_{p\le y}(\log p)/p\ll\log y$ and
Equation~\eqref{eq:sieve-local-tail} give a fixed $C_*\ge1$ such that
\[
 \sum_{p\in\mathcal P_y}\vartheta_p\log p\le C_*\log y
 \qquad(y\ge2).
\]
Choose fixed positive parameters satisfying
\begin{equation}\label{eq:sieve-parameters}
 20(C+1)\eta<\theta,\qquad
 0<\delta<\min\{\eta/4,\eta/(4C_*)\},
 \qquad T_0=X^\eta,\quad Y=X^\delta.
\end{equation}
Here $T_0$ truncates the squarefree Selberg divisors and is the stopping
threshold for depth products, whose crossing prefixes may reach $T_0^2$;
$Y$ bounds the small primes outside $S$ handled by the sieve.
For every $\mathcal I\subset\mathcal P_Y$, Mertens' estimate and
$-\log(1-\vartheta_p)\le2\vartheta_p\le2C/p$ give, for a fixed $C_1$,
\begin{equation}\label{eq:sieve-inactivity-lower}
 \prod_{p\in\mathcal I}(1-\vartheta_p)
       \gg(\log X)^{-C_1}.
\end{equation}
The bound is uniform in the subset $\mathcal I$ and in $S$.

\begin{lemma}[Prescribed depths and inactive primes]\label{lem:sieve-assignment}
Let $\mathcal A,\mathcal I$ be disjoint subsets of $\mathcal P_Y$, and
assign an integer $j_p\ge1$ to each $p\in\mathcal A$.  Suppose that
$\prod_{p\in\mathcal A}p^{j_p}\le T_0^2$.  Then
\begin{equation}\label{eq:sieve-assignment}
 \int
   \prod_{p\in\mathcal A}
      b_p\mathbf 1_{\{n_p=j_p\}}
   \prod_{p\in\mathcal I}\mathbf 1_{\{n_p=0\}}\,
        d\nu_\phi
 \ll V_\phi
   \prod_{p\in\mathcal A}\E_p
      \bigl(b_p\mathbf 1_{\{n_p=j_p\}}\bigr)
   \prod_{p\in\mathcal I}(1-\vartheta_p).
\end{equation}
The constant is uniform in the two subsets and all their assigned depths.
\end{lemma}

\begin{proof}
We use the Selberg upper-bound square
\cite{Selberg1947}; see also \cite[Chapter~21]{Koukoulopoulos2019}.
Its calculation is included because the relative error in
Lemma~\ref{lem:sieve-residue-function} must preserve the assigned means.
If $\vartheta_p=0$, the activity event at $p$ is a union of residue classes
of Haar measure zero, hence is empty on global points.  Such a prime can
be removed from $\mathcal I$.  For the remaining primes put
\[
 \vartheta(d)=\prod_{p\mid d}\vartheta_p,\qquad
 h(d)=\prod_{p\mid d}\frac{\vartheta_p}{1-\vartheta_p}
\]
for squarefree $d$ supported on $\mathcal I$.  Both functions take value
one at $d=1$.  Let $\mathcal D$ be the set of those integers with
$d\le T_0$, and put $\mathcal H=\sum_{d\in\mathcal D}h(d)$.

For coefficients $\lambda_d$ supported on $\mathcal D$ with
$\lambda_1=1$, the square
\[
 \left(\sum_{d\in\mathcal D}\lambda_d
       \prod_{p\mid d}\mathbf 1_{\{n_p\ge1\}}\right)^2
\]
majorizes the indicator that every prime of $\mathcal I$ is inactive.
Its expectation in the independent product of the local probability
spaces is
\begin{equation}\label{eq:sieve-quadratic}
 \sum_{d,e\in\mathcal D}\lambda_d\lambda_e\vartheta([d,e])
   =\sum_{l\in\mathcal D}\frac1{h(l)}
       \left(\sum_{\substack{d\in\mathcal D\\l\mid d}}
                      \lambda_d\vartheta(d)\right)^2.
\end{equation}
Here $[d,e]$ denotes the least common multiple.  The identity follows
by expanding the right side and using
$\vartheta([d,e])
 =\vartheta(d)\vartheta(e)\sum_{l\mid(d,e)}h(l)^{-1}$.

The set $\mathcal D$ is closed under divisors.  Finite M\"obius inversion
therefore permits the choice
\[
 \sum_{\substack{d\in\mathcal D\\l\mid d}}\lambda_d\vartheta(d)
       =\frac{\mu_{\mathrm{Mob}}(l)h(l)}{\mathcal H},
 \qquad l\in\mathcal D,
\]
where $\mu_{\mathrm{Mob}}$ is the M\"obius function.  Explicitly,
\[
 \lambda_d=
  \frac{\mu_{\mathrm{Mob}}(d)h(d)}
       {\vartheta(d)\mathcal H}
  \sum_{\substack{r\le T_0/d\\(r,d)=1\\dr\in\mathcal D}}h(r).
\]
Consequently $\lambda_1=1$, the value of
Equation~\eqref{eq:sieve-quadratic} is $\mathcal H^{-1}$, and
\begin{equation}\label{eq:sieve-coefficients}
 |\lambda_d|
 \le\frac{h(d)}{\vartheta(d)}
 =\prod_{p\mid d}(1-\vartheta_p)^{-1}
 \le2^{\omega(d)}\le T_0.
\end{equation}

To estimate $\mathcal H$, give every squarefree integer supported on
$\mathcal I$ probability proportional to $h(d)$, without the bound
$d\le T_0$.  In this finite product distribution, prime $p$ is selected
with probability $\vartheta_p$.  Thus
\[
 \E\log d=\sum_{p\in\mathcal I}\vartheta_p\log p
       \le C_*\log Y<\tfrac14\log T_0.
\]
Markov's inequality shows that at least half the mass lies at $d\le T_0$.
It follows that
\begin{equation}\label{eq:sieve-selberg-main}
 \mathcal H\ge\tfrac12\prod_{p\in\mathcal I}(1+h(p)),
 \qquad
 \mathcal H^{-1}\le2\prod_{p\in\mathcal I}(1-\vartheta_p).
\end{equation}

It remains to average this square against $\nu_\phi$ together with the
assigned weights.  Write
\[
 A_{\mathcal A}=
   \prod_{p\in\mathcal A}\E_p
        \bigl(b_p\mathbf 1_{\{n_p=j_p\}}\bigr).
\]
The assigned weight is defined modulo
\[
 \prod_{p\in\mathcal A}p^{C(j_p+1)}
 \le\left(\prod_{p\in\mathcal A}p^{j_p}\right)^{2C}
 \le T_0^{4C}.
\]
For a pair $d,e\in\mathcal D$, its activity indicators require only the
additional modulus $[d,e]^{2C}\le T_0^{4C}$.  The two prime sets are
disjoint, so the total modulus is at most
\begin{equation}\label{eq:sieve-modulus-budget}
 T_0^{8C}\le X^\theta
\end{equation}
and is prime to every prime of $S$.  The function before multiplication
by $\lambda_d\lambda_e$ is nonnegative, and its residue mean is
$A_{\mathcal A}\vartheta([d,e])$ by the Chinese remainder theorem.
Lemma~\ref{lem:sieve-residue-function} therefore gives the corresponding
main term $V_\phi A_{\mathcal A}\vartheta([d,e])$ and an error of absolute
value at most
$C V_\phi X^{-\theta}A_{\mathcal A}\vartheta([d,e])$.
There are at most $T_0^2$ pairs and, by
Equation~\eqref{eq:sieve-coefficients}, each coefficient product has
absolute value at most $T_0^2$.  Summing the errors absolutely costs at
most
\begin{equation}\label{eq:sieve-square-error}
 C V_\phi A_{\mathcal A}X^{-\theta}T_0^4.
\end{equation}
This estimate keeps $A_{\mathcal A}$ even if it is very small.
The main term is $V_\phi A_{\mathcal A}\mathcal H^{-1}$ by
Equation~\eqref{eq:sieve-quadratic}.  Finally,
$X^{-\theta}T_0^4=o((\log X)^{-C_1})$ by
Equation~\eqref{eq:sieve-parameters}.  Equations
\eqref{eq:sieve-inactivity-lower} and
\eqref{eq:sieve-selberg-main} absorb
Equation~\eqref{eq:sieve-square-error} and prove the lemma.
\end{proof}

\subsection{Truncating and summing the depths}

\begin{proof}[Proof of Proposition~\ref{prop:sieve-averaging}]
Put
$P_Y=\prod_{p\in\mathcal P_Y}\E_p b_p$.  Equations
\eqref{eq:sieve-mean-lower} and
\eqref{eq:sieve-inactivity-lower} imply
\begin{equation}\label{eq:sieve-product-lower}
 P_Y\gg(\log X)^{-C_1}.
\end{equation}
First observe that primes exceeding $Y$ cost only a fixed factor.
Equations~\eqref{eq:sieve-depth-product} and
\eqref{eq:sieve-pointwise-weight}, together with $b_p\le A_0^{n_p}$, give
\begin{equation}\label{eq:sieve-large-primes}
 \prod_{\substack{p>Y\\p\notin S}}b_p(z)
 \le A_0^{\sum_{p>Y,\ p\notin S}n_p(z)}
 \le \exp(B\log A_0/\delta).
\end{equation}

Consider the points for which
$\prod_{p\in\mathcal P_Y}p^{n_p(z)}\le T_0$.
Assign all their active primes in $\mathcal P_Y$ and require inactivity
at the other primes in $\mathcal P_Y$.  These assignments satisfy
Lemma~\ref{lem:sieve-assignment}.  After applying it, we may enlarge the
sum of its nonnegative upper bounds by allowing every subset and every
positive depth.  The resulting product is
\[
 V_\phi\prod_{p\in\mathcal P_Y}
   \left(1-\vartheta_p+
       \sum_{j\ge1}\E_p(b_p\mathbf 1_{\{n_p=j\}})\right)
   =V_\phi P_Y.
\]
Together with Equation~\eqref{eq:sieve-large-primes}, this bounds the
contribution from these points.

For a remaining point, the product through $\mathcal P_Y$ exceeds $T_0$.
Order its active primes in $\mathcal P_Y$ increasingly and stop at the
first prefix whose product exceeds $T_0$.  If every individual factor
$p^{n_p(z)}$ for $p\in\mathcal P_Y$ is at most $T_0$, the prefix has product in
$(T_0,T_0^2]$, since its preceding product and its last factor are both
at most $T_0$.  If also
$\prod_{p\in\mathcal P_{y_0}}p^{n_p(z)}\le T_0$, for a fixed cutoff
$y_0>4$ to be chosen below independently of $X$, its last prime exceeds
$y_0$.  The next two estimates allow these restrictions at negligible
weighted cost.

First, if $p^{n_p(z)}>T_0$ for $p\in\mathcal P_Y$, then
$n_p(z)\ge J_p$, where
$J_p=\lfloor\log T_0/\log p\rfloor+1$.  The threshold is determined modulo
\[
 p^{C(J_p+1)}\le T_0^C Y^{2C}\le X^\theta.
\]
Lemma~\ref{lem:sieve-residue-function} and
Equation~\eqref{eq:sieve-local-tail} bound its $\nu_\phi$-measure by
$O(V_\phi p^{-J_p})=O(V_\phi/T_0)$.  Summing over $p\in\mathcal P_Y$
gives $O(V_\phi Y/T_0)$.  By Equation~\eqref{eq:sieve-pointwise-weight},
the contribution after multiplication by $\mathcal W$ is
$O_\epsilon(V_\phi X^{\epsilon+\delta-\eta})$.
Choose $\epsilon<\eta-\delta$.  This is smaller than $V_\phi P_Y$ by a
power of $X$, in view of Equation~\eqref{eq:sieve-product-lower}.

We also need the analogous removal for a fixed nonempty finite set
$\mathcal Q$ of primes, even though some of them may belong to the varying
set $S$.  Put $N=|\mathcal Q|$.  If
$\prod_{p\in\mathcal Q\setminus S}p^{n_p(z)}>T_0$, some prime in
$\mathcal Q\setminus S$ has $p^{n_p(z)}>T_0^{1/N}$.  Put
$J_{p,\mathcal Q}=\lfloor\log T_0/(N\log p)\rfloor+1$.  The event
$n_p\ge J_{p,\mathcal Q}$ is defined modulo at most
$T_0^{C/N}p^{2C}\le X^\theta$ for large $X$, and its local probability
is at most $Cp^{-J_{p,\mathcal Q}}\le C T_0^{-1/N}$.
Lemma~\ref{lem:sieve-residue-function} therefore bounds the measure of
the union by
$O_{\mathcal Q}(V_\phi T_0^{-1/N})$.  After multiplication by
Equation~\eqref{eq:sieve-pointwise-weight} with $\epsilon<\eta/N$, its
weighted contribution again saves a power of $X$ relative to $V_\phi$.
If $\mathcal Q\setminus S$ is empty the event is empty.  Thus this
removal is uniform in the possible intersections of $\mathcal Q$ with $S$.

Choose a fixed $M_0>0$ and then a fixed $y_0>4$ such that
\begin{equation}\label{eq:sieve-crossing-parameters}
 \eta M_0>2B\log A_0+1,\qquad M_0/\log y_0<\tfrac14.
\end{equation}
Apply the finite-set removal with $\mathcal Q=\{p:p\le y_0\}$.
These choices precede the choice of the small $\epsilon$ used in that
removal, and all are independent of $X$.

The retained first-crossing prefixes therefore have product in
$(T_0,T_0^2]$ and last prime greater than $y_0$.  Group these assignments
by the interval containing their last prime:
\begin{equation}\label{eq:sieve-buckets}
 \sqrt y<p\le y,\qquad y=Y^{1/2^i},\quad i=0,1,2,\ldots.
\end{equation}
Only groups with $y>y_0$ occur.  Every unassigned active prime is greater
than the last assigned prime and hence greater than $\sqrt y$.  The total
weight of those primes, including primes beyond $Y$, is therefore at most
\begin{equation}\label{eq:sieve-unassigned}
 \exp\left(\frac{2B\log A_0}{\log y}\log X\right).
\end{equation}
For each assignment in the group, require inactivity at its unassigned
primes in $\mathcal P_{\sqrt y}$ and apply
Lemma~\ref{lem:sieve-assignment}.  We use Rankin's power majorant
\cite{Rankin1938}: since the assigned product exceeds $T_0$, its upper
bound can be multiplied by
\[
 T_0^{-s}\prod_{\text{assigned }p}p^{s j_p}\ge1,
 \qquad s=M_0/\log y.
\]
Now enlarge the sum of these nonnegative upper bounds to all subsets and
all positive depths through $y$.  This enlarges only a numerical sum after
the lemma has been applied to the genuine prefixes; it invokes no residue
estimate for the newly admitted products above $T_0^2$.

For $p\in\mathcal P_{\sqrt y}$, the option of no assignment has weight
$1-\vartheta_p$ and the assigned options have total weight
$\sum_{j\ge1}\E_p(b_p\mathbf 1_{\{n_p=j\}})p^{sj}$.  Their sum is
$\E_p(b_p p^{s n_p})$.  For $\sqrt y<p\le y$ outside $S$, inactivity was
not imposed, so the option of no assignment has weight one.  The total
there is $\E_p(b_p p^{s n_p})+\vartheta_p$.  Thus the group's contribution
is at most
\begin{equation}\label{eq:sieve-bucket-bound}
 C V_\phi T_0^{-s}
 \exp\left(\frac{2B\log A_0}{\log y}\log X\right)
 \prod_{p\in\mathcal P_{\sqrt y}}\E_p(b_p p^{s n_p})
 \prod_{\substack{\sqrt y<p\le y\\p\notin S}}
       \bigl(\E_p(b_p p^{s n_p})+\vartheta_p\bigr).
\end{equation}

We estimate these local moments uniformly for $y>y_0$.  By
Equation~\eqref{eq:sieve-crossing-parameters}, $s<1/4$.  The total
contribution of depths $j\ge2$ to $\E_p(b_p p^{s n_p})$ is at most
\[
 C^2\sum_{j\ge2}(1+j)^C p^{-(1-s)j}\ll_C p^{-3/2}.
\]
At depth one, the increase over $\E_p b_p$ is bounded by
\[
 \frac{C'}p(p^s-1)
   \ll_{M_0}\frac{\log p}{p\log y}
 \qquad(p\le y),
\]
because $0\le s\log p\le M_0$.  Chebyshev's estimate shows that these
increases have bounded sum over $p\le y$.  Mertens' estimate also gives
$\sum_{\substack{\sqrt y<p\le y\\p\notin S}}\vartheta_p\ll1$.  Since
$\E_p b_p\ge1/2$, comparison of the products in
Equation~\eqref{eq:sieve-bucket-bound} with their base means gives
\[
 \prod_{p\in\mathcal P_{\sqrt y}}\E_p(b_p p^{s n_p})
 \prod_{\substack{\sqrt y<p\le y\\p\notin S}}
       \bigl(\E_p(b_p p^{s n_p})+\vartheta_p\bigr)
 \ll_{M_0}\prod_{p\in\mathcal P_y}\E_p b_p.
\]
For a fixed $C_2$, a final use of
Equation~\eqref{eq:sieve-mean-lower} and Mertens' estimate yields
\[
 \prod_{p\in\mathcal P_y}\E_p b_p
 \le P_Y\prod_{\substack{y<p\le Y\\p\notin S}}
                 (1-\vartheta_p)^{-1}
 \ll P_Y\left(\frac{\log Y}{\log y}\right)^{C_2}.
\]
By Equation~\eqref{eq:sieve-crossing-parameters}, the bound for the group
in Equation~\eqref{eq:sieve-bucket-bound} is consequently at most
\[
 C_{M_0}V_\phi P_Y
 \exp\left(-\frac{\log X}{\log y}\right)
 \left(\frac{\log Y}{\log y}\right)^{C_2}.
\]
For the values of $y$ in Equation~\eqref{eq:sieve-buckets}, its sum is
bounded by
\[
 C_{M_0}V_\phi P_Y
       \sum_{i\ge0}\exp(-2^i/\delta)\,2^{iC_2}
 \ll V_\phi P_Y.
\]
Together with the small-product part and the two discarded sets, this
proves Equation~\eqref{eq:sieve-average}.
\end{proof}

\subsection{The planar exterior count}

\begin{proof}[Proof of Theorem~\ref{thm:exterior-count}]
Fix the nonzero trace and bounded region in that theorem.  Let
$\mathcal R$ be any of its rectangles, in the specified two partition
coordinates.  Choose an allowed smooth function $\phi$ that majorizes
$\mathbf 1_{\mathcal R}$, with its two scales equal to the side lengths
of $\mathcal R$.  Its support lies in a fixed enlargement of the bounded
region and
\[
 A_\phi\asymp\operatorname{area}(\mathcal R).
\]
The constants are uniform over the positions, side lengths, and aspect
ratios of the rectangles.  Proposition~\ref{prop:ext-multiplicity-majorant}
and Proposition~\ref{prop:sieve-averaging} give
\begin{equation}\label{eq:sieve-before-euler}
 \sum_{\substack{\Tr z=m\\z\in\mathcal R}}M(z)
 \ll \sqrt D\, A_\phi
       \prod_{p\in S}(1-p^{-1})Z_p
       \prod_{p\in\mathcal P_Y}\E_p b_p,
 \qquad
 Z_p=\prod_{\substack{\mathfrak p\mid p\\\mathfrak p\subset R}}
                (1-\Nm(\mathfrak p)^{-1})^{-1}.
\end{equation}
Here $\prod_{p\in S}|D|_p^{-1/2}=\sqrt D$, because $S$ contains every
prime divisor of $D$, and Equation~\eqref{eq:sieve-volume} absorbed the
trace densities.  The restriction to integral points outside $S$ in the
multiplicity proposition agrees with the support of $\nu_\phi$.

The exact unramified mean in Lemma~\ref{lem:ext-unramified-data} is
\[
 \E_p b_p=(1-p^{-1})(1-p^{-3})Z_p\qquad(p\notin S).
\]
Since $1-p^{-3}\le1$, the two Euler products on the right of
Equation~\eqref{eq:sieve-before-euler} are at most the product of
$(1-p^{-1})Z_p$ over $p\in S\cup\mathcal P_Y$.
Ideal factorization in the degree-four maximal order gives, uniformly
also at ramified primes,
\[
 \log Z_p=\frac{a_K(p)}p+O(p^{-2}).
\]
Thus
\begin{equation}\label{eq:sieve-euler}
 \prod_{p\in S}(1-p^{-1})Z_p
       \prod_{p\in\mathcal P_Y}\E_p b_p
 \ll
 \exp\left(\sum_{p\in S\ \mathrm{or}\ p\le Y}
                    \frac{a_K(p)-1}{p}\right).
\end{equation}

For fixed $m$, the radical of $S$ is polynomial in $X$.  Choose a fixed
$c\ge1/2$ large enough that, for all sufficiently large $X$, every prime
in $S$ and every prime at most $Y$ is at most $Z=X^c$.  Since
$a_K(p)\ge0$, the exponent in Equation~\eqref{eq:sieve-euler} is at most
\[
 \sum_{p\le Z}\frac{a_K(p)}p-\sum_{p\le Y}\frac1p.
\]
We have $Z\ge D^{1/4}$ because $X=q\sqrt D\ge\sqrt D$.
Equation~\eqref{eq:ord-prime-sum} and Mertens' estimate now give
\[
 \exp\left(\sum_{p\in S\ \mathrm{or}\ p\le Y}
                    \frac{a_K(p)-1}{p}\right)
 \ll \frac{\kappa\log Z}{\log Y}
 =\frac c\delta\,\kappa.
\]
This also applies when $D$ remains bounded and $q$ tends to infinity,
by the stated uniform range of Lemma~\ref{lem:ord-residue}.
Substitution in Equation~\eqref{eq:sieve-before-euler} proves
\[
 \sum_{\substack{\Tr z=m\\z\in\mathcal R}}M(z)
       \ll\kappa\sqrt D\,\operatorname{area}(\mathcal R).
\]
The only box-dependent quantity in the proof is $A_\phi$, which is
comparable to the area with a fixed constant.  All other estimates are
uniform over the box translations and the two scales between
$X^{-\tau}$ and the fixed upper bounds.  The local multiplicity majorant
has the same constant for every chosen ideal representative system, and
its right side is independent of that choice.  Thus the bound is uniform
in the representative system as well.  This proves the precise
uniformity asserted in Theorem~\ref{thm:exterior-count}.
\end{proof}
